\chapter{Revisão da Literatura}
A revisão da literatura estabelece as bases teóricas e conceituais para a integração de DevSecOps e Infraestrutura como Código (IaC), abordando a evolução dessas práticas e sua relevância no contexto de automação e segurança. Este capítulo explora conceitos fundamentais, como o papel da segurança no ciclo de vida do desenvolvimento de software (DevSecOps) e a gestão automatizada de infraestrutura baseada em código (IaC). Por fim, destaca-se a convergência dessas abordagens como uma solução abrangente para os desafios contemporâneos em TI, oferecendo consistência, eficiência e governança em ambientes computacionais complexos.

\begin{enumerate}
    \item\textbf{DevSecOps:} O conceito de DevSecOps visa incorporar a segurança como um aspecto essencial e integrado ao ciclo de vida de desenvolvimento de software (SDLC). Isso envolve automação, colaboração entre equipes e adoção de princípios como \textit{Shift-left}, onde a segurança é abordada desde as fases iniciais do desenvolvimento;
    \item\textbf{Infraestrutura como Código (IaC):} IaC é a prática de gerenciar a infraestrutura de TI por meio de arquivos de configuração legíveis por humanos, permitindo maior consistência, reprodutibilidade e controle de mudanças;
    \item\textbf{Integração de DevSecOps e IaC:} A integração dessas abordagens oferece uma solução abrangente para desafios de automação e segurança em ambientes computacionais complexos;
\end{enumerate}

A seguir, detalha-se cada uma dessas dimensões com base na literatura existente.



\section{DevSecOps}
A evolução do DevOps, que unifica desenvolvimento e operações para entregar software de maneira mais rápida e confiável, levou à necessidade de integrar segurança ao ciclo de vida do desenvolvimento.

As principais práticas incluem:

\begin{itemize}
    \item\textbf{\textit{Shift-left:}} Priorização de testes e revisões de segurança nas fases iniciais do ciclo de vida.
    \item\textbf{Automatização:} Uso de ferramentas como scanners de vulnerabilidades e integração de segurança nos pipelines CI/CD.
    \item\textbf{Colaboração:} Promoção de uma cultura onde a segurança é responsabilidade de todos.
\end{itemize}

DevSecOps enfatiza a incorporação da segurança desde o início do ciclo de vida de desenvolvimento, utilizando automação e colaboração contínuas. \citeonline{kim2021devops}, destacam que a segurança deve ser tratada como um problema compartilhado entre as equipes de desenvolvimento e operações, com práticas integradas diretamente no fluxo de trabalho diário. Adicionar segurança ao DevOps requer mais do que ferramentas; requer uma mudança cultural que priorize a segurança sem comprometer a velocidade e a agilidade\cite{vehent2018securing}.

\subsection{Shift-left}
O conceito de Shift Left refere-se à prática de integrar atividades de segurança e qualidade no início do ciclo de desenvolvimento de software, ao invés de tratá-las como etapas finais. Essa abordagem busca identificar e corrigir problemas de segurança durante as fases de planejamento, design e codificação, o que reduz custos e aumenta a eficiência. Para isso, ferramentas automatizadas como scanners de código estático (SAST) e análise de dependências são usadas para detectar vulnerabilidades antes que o software avance para etapas mais complexas, como testes ou implantação.

Essa mudança para a esquerda no pipeline de desenvolvimento promove uma cultura de "segurança como responsabilidade compartilhada", envolvendo desenvolvedores, equipes de qualidade e operações desde o início. Além disso, práticas como testes unitários com foco em segurança, integração contínua com validações de segurança (CI), e treinamento dos desenvolvedores sobre vulnerabilidades comuns, como as listadas no OWASP, são pilares do \textit{Shift Left}. Ao adotar essa mentalidade, as organizações conseguem entregar produtos mais seguros e com maior qualidade, alinhando-se aos princípios fundamentais do DevSecOps.



\section{Infraestrutura como Código (IaC)}
IaC transforma a gestão da infraestrutura, substituindo processos manuais por automação baseada em arquivos de configuração versionados, o que contribui para:

\begin{itemize}
    \item\textbf{Reprodutibilidade:} Reduzindo erros humanos e permitindo que sistemas sejam provisionados de forma idêntica em diferentes ambientes.
    \item\textbf{Gestão de configurações:} Garantindo consistência e controle de mudanças.
    \item\textbf{Segurança:} Incorporando políticas de segurança diretamente nos códigos de configuração.
\end{itemize}

As principais ferramentas de IaC incluem Terraform, Ansible e CloudFormation, cada uma com características específicas que atendem a diferentes necessidades organizacionais.

IaC permite gerenciar configurações e recursos de infraestrutura de forma declarativa, garantindo consistência e facilidade de reprodução. Segundo \citeonline{morris2016infrastructure} infraestrutura tratada como código é fundamental para habilitar mudanças rápidas, seguras e confiáveis, permitindo que as equipes colaborem e iterem continuamente. Além disso, IaC reduz erros humanos e fornece documentação implícita ao transformar a infraestrutura em algo que pode ser versionado e testado como qualquer outro código\cite{brikman2022terraform}.


\section{Integração de DevSecOps e IaC}
A combinação de DevSecOps e IaC permite um modelo em que a segurança é implementada automaticamente durante o provisionamento e a gestão de infraestrutura.

\begin{itemize}
    \item\textbf{Detecção precoce de vulnerabilidades:} Integração de scanners de segurança em pipelines CI/CD.
    \item\textbf{Conformidade automatizada:} Aderência a normas e regulamentações por meio de políticas incorporadas.
    \item\textbf{Eficiência operacional:} Redução de custos e tempo associado a falhas de segurança.
\end{itemize}

Ferramentas como Checkov e Open Policy Agent (OPA) auxiliam na verificação de conformidade e na aplicação de políticas de segurança em códigos IaC.

 A união de DevSecOps e IaC promete abordar questões críticas de segurança e automação, possibilitando melhorias significativas em operações de TI. \citeonline{domingus2022cloud} afirmam que a convergência de DevOps, segurança e infraestrutura como código cria um ciclo de feedback contínuo, onde a automação e a governança trabalham em harmonia. Em linha com isso, \citeonline{forsgren2018accelerate} enfatizam que as organizações que integram segurança no fluxo de trabalho e usam práticas modernas, como IaC, apresentam um desempenho significativamente superior.